\newcommand{\DoNotLoadEpstopdf}{}
\documentclass[11pt]{article}
\usepackage[T1]{fontenc}
\usepackage{lmodern}
\usepackage[margin=1.08in]{geometry}
\usepackage{amsmath,amssymb,amsthm,mathtools}
\usepackage{microtype}
\usepackage{xurl}
\usepackage{xcolor}
\usepackage[colorlinks=true,linkcolor=blue!50!black,citecolor=blue!50!black,urlcolor=blue!50!black]{hyperref}
\input{glyphtounicode}
\input{glyphtounicode-cmex}
\pdfgentounicode=1
\pdftrailerid{}
\pdfsuppressptexinfo=15
\hypersetup{pdftitle={An L2 Einstein Gap for Nonnegatively Curved Four-Manifolds},pdfauthor={OpenAI}}
\newtheorem{theorem}{Theorem}[section]
\newtheorem{proposition}[theorem]{Proposition}
\newtheorem{lemma}[theorem]{Lemma}
\theoremstyle{definition}
\newtheorem{assumption}[theorem]{Classification Assumption}
\theoremstyle{remark}
\newtheorem{remark}[theorem]{Remark}
\numberwithin{equation}{section}
\DeclareMathOperator{\Ric}{Ric}
\DeclareMathOperator{\Rm}{Rm}
\DeclareMathOperator{\Scal}{Scal}
\DeclareMathOperator{\Vol}{Vol}
\DeclareMathOperator{\diam}{diam}
\DeclareMathOperator{\inj}{inj}
\DeclareMathOperator{\diver}{div}
\newcommand{\dd}{\,d\mu}
\newcommand{\R}{\mathbb R}
\newcommand{\CP}{\mathbb{CP}}
\newcommand{\E}{\mathcal E}
\newcommand{\norm}[1]{\left\lVert #1\right\rVert}
\title{An $L^2$ Einstein Gap for\\Nonnegatively Curved Four-Manifolds}
\author{OpenAI}
\date{October 5, 2026}

\begin{document}
\maketitle
\begin{abstract}
We prove a universal, scale-invariant $L^2$ gap for the trace-free Ricci
tensor on simply connected closed four-manifolds with nonnegative
sectional curvature. If the trace-free Ricci energy is sufficiently
small, the manifold is diffeomorphic to $S^4$, $\CP^2$, or
$S^2\times S^2$. The proof uses the classification of positive-Einstein,
nonnegatively curved four-manifolds supplied by the companion
zero-plane rigidity theorem, stated explicitly as the classification
premise of our result. No auxiliary curvature, volume, diameter,
injectivity-radius, or Sobolev bound is required. The conclusion concerns
the smooth manifold, not an isometry of the original metric.
\end{abstract}

\section{Introduction}
\label{sec:intro}

For a Riemannian four-manifold $(M,g)$, let
\[
 E_g=\Ric_g-\frac{\Scal_g}{4}g,
 \qquad \mathcal D(g)=\int_M |E_g|_g^2\,d\mu_g.
\]
The tensor $E_g$ vanishes precisely when $g$ is Einstein, provided that
$M$ is connected. The energy $\mathcal D$ is unchanged by constant
rescaling in dimension four. It is therefore natural to ask whether a
classification of Einstein metrics remains valid at the level of smooth
manifolds when $\mathcal D$ is small. We address this question under
nonnegative sectional curvature.

We take the exact Einstein classification in
\cite[Corollary~1.2]{ZeroPlane} as the following premise.

\begin{assumption}\label{ass:classification}
Every connected smooth closed Riemannian four-manifold $(N,h)$ satisfying
\[
 \Ric_h=3h,\qquad \sec_h\ge0
\]
has universal Riemannian cover isometric to one of the following:
the round $S^4$ of sectional curvature $1$; Fubini--Study $\CP^2$ normalized
by $\Ric=3g$; or
$S^2(1/\sqrt3)\times S^2(1/\sqrt3)$ with its product metric.
\end{assumption}

This premise is used only for an exactly Einstein metric; no stability
statement is assumed.

\begin{theorem}\label{thm:gap}
Assume Classification Assumption~\ref{ass:classification}. There is a
universal constant $\varepsilon_0>0$ with the following property.
If $(M,g)$ is a connected, simply connected, smooth closed Riemannian
four-manifold with $\sec_g\ge0$ and
\[
 \int_M\left|\Ric_g-\frac{\Scal_g}{4}g\right|_g^2\,d\mu_g
 <\varepsilon_0,
\]
then $M$ is diffeomorphic, without a prescribed orientation, to the
standard smooth $S^4$, $\CP^2$, or $S^2\times S^2$.
\end{theorem}

There are no auxiliary bounds on curvature, volume, diameter, injectivity
radius, Sobolev constants, or curvature derivatives. The conclusion is
about the underlying smooth manifold, rather than the original metric.
This distinction matters because small $L^2$ error does not itself give
pointwise curvature control.

\subsection*{Context and methods}

The exact Einstein problem has been studied under several additional
hypotheses. Gursky and LeBrun obtained the Fubini--Study conclusion for
nonnegative sectional curvature and a nonzero positive-definite
intersection form \cite{GurskyLeBrun}. Cao and Tran established
sectional-pinching criteria \cite{CaoTran}, and Liu classified the
$T^2$-invariant case on closed simply connected four-manifolds
\cite{Liu}. Here the classification premise is supplied by
\emph{Zero-Plane Rigidity for Einstein Four-Manifolds}
\cite[Corollary~1.2]{ZeroPlane}. Our problem is instead to transfer
that exact classification to metrics with small critical Ricci error.

The relation between Ricci curvature, scalar curvature, and topology is
especially strong in dimension four: the Chern--Gauss--Bonnet formula
controls the Euler characteristic by the total squared curvature.
The almost-Schur inequality of De Lellis and Topping
\cite{DeLellisTopping} controls the variance of scalar curvature by the
trace-free Ricci tensor under nonnegative Ricci curvature. We use its
Bianchi--Bochner argument to choose a scale in which the Ricci tensor is
close in $L^2$ to $3g$.

Integral versions of Myers' diameter argument combine second variation
with control of curvature along many geodesics. The segment inequality
of Cheeger and Colding \cite{CheegerColding} and the integral-curvature
diameter estimates of Petersen and Sprouse \cite{PetersenSprouse,Sprouse} provide
the relevant context. We give the short segment-averaging argument needed
here, together with an elementary derivation of the initial Sobolev
inequality. These estimates prevent collapse before any smoothing is
performed.

Hamilton's Ricci flow \cite{HamiltonRF} supplies the regularization.
Perelman's entropy monotonicity \cite{Perelman} propagates a
scalar-curvature-weighted Sobolev inequality. The passage from entropy
to Sobolev inequalities along Ricci flow is developed by Ye and Zhang
\cite{Ye,Zhang}; we include the form needed here, with constants that are
uniform even before a common existence interval is known. Shi's
derivative estimates \cite{Shi} and Hamilton's compactness theorem
\cite{HamiltonCompactness} then turn curvature control at positive times
into smooth convergence.

Smoothing from small critical curvature concentration was developed,
in particular, by Chan, Chen, and Lee \cite{ChanChenLee}. The formulation
used here isolates a globally small \emph{high part} of the curvature.
Consider a sequence $(M_j,g_j)$ satisfying the geometric hypotheses of
Theorem~\ref{thm:gap} with $\mathcal D(g_j)\to0$. After normalization,
nonnegative sectional curvature gives, for a fixed threshold $K_0$,
\[
 \norm{(|\Rm_{g_j}|-K_0)_+}_{L^2(g_j)}\longrightarrow0.
\]
A moving threshold propagates this estimate. It yields both a uniform
Sobolev inequality and noncollapse along each existing flow; a
curvature-concentration contradiction then establishes a common
lifespan. The same tail bound controls the evolving Einstein error.
Thus the high-part formulation supplies both the regularization and
the defect estimate needed for the gap argument.

A second point concerns sectional curvature. Its nonnegativity is not
assumed to persist under the flow. Instead, we bound the integral of its
negative part by $Ct$. The limit is an exactly Einstein shrinking flow,
and its explicit scaling converts this weak defect estimate into
nonnegative sectional curvature of the fixed Einstein metric. This gives
the classification-independent compactness statement of
Proposition~\ref{prop:limit}; classification then identifies the limiting
manifold.

\subsection*{Proof organization}

Section~\ref{sec:initial} normalizes the sequence and establishes the
initial geometric and Sobolev bounds. Section~\ref{sec:smoothing} obtains
a uniform smoothing interval by controlling the high-curvature part.
Section~\ref{sec:defects} propagates the Einstein error and estimates the
negative sectional curvature. Section~\ref{sec:limit} constructs the
compact Einstein limit, identifies its sectional curvature, and proves
Theorem~\ref{thm:gap}.

All manifolds are smooth, and \emph{closed} means compact without
boundary. We use $R=\Scal$ and the curvature convention in which
$\Rm(X,Y,X,Y)$ is the sectional curvature of an orthonormal pair.
The Laplacian is $\Delta=\operatorname{tr}\nabla^2$. Norms and integrals
use the displayed metric, or the current metric along a flow, and are
not volume-normalized. Constants denoted by $C$ or $c$ may change from
line to line. Uniformity always means independence of the sequence
index, after discarding finitely many terms.

\section{Normalization and initial geometry}
\label{sec:initial}

Throughout the proof let $(M_j,g_j)$ be connected, simply connected,
closed four-manifolds with
\begin{equation}\label{eq:sequence}
 \sec_{g_j}\ge0,\qquad \norm{E_{g_j}}_2\longrightarrow0.
\end{equation}
The aim of this section is to obtain uniform geometric bounds without
assuming a curvature maximum bound. We suppress $j$ when considering
one metric in the sequence.

\begin{proposition}\label{prop:initial}
After constant rescaling, the metrics in \eqref{eq:sequence} satisfy
\begin{equation}\label{eq:normalization}
 \frac1V\int_M R\dd=12,\qquad
 \norm{R-12}_2\le4\norm E_2,\qquad
 \norm{\Ric-3g}_2\longrightarrow0,
\end{equation}
where $V=\Vol(M,g)$. There are uniform positive constants
$v_0,V_0,D_0$ such that
\begin{equation}\label{eq:initial-geometry}
 v_0\le V\le V_0,\qquad \diam(M,g)\le D_0.
\end{equation}
\end{proposition}

\begin{proof}
Simple connectivity implies orientability and $b_1=0$. Poincar\'e duality
therefore gives $\chi(M)=2+b_2(M)\ge2$. By Chern--Gauss--Bonnet,
\begin{equation}\label{eq:euler}
 2\le\chi(M)\le C\int_M|\Rm|^2\dd.
\end{equation}
On the space of algebraic curvature tensors, the maximum absolute
sectional curvature is a norm: sectional curvatures determine the tensor.
Equivalence of norms in this finite-dimensional space, together with
$0\le\sec\le R/2$, gives a dimensional constant $C_0$ such that
\begin{equation}\label{eq:algebraic}
 |\Rm|\le C_0R.
\end{equation}
The average scalar curvature is positive, since otherwise $R\equiv0$
and \eqref{eq:algebraic} contradicts \eqref{eq:euler}. Rescale so that
its average is $12$. The $L^2$ norm of $E$ is unchanged.

For completeness, the almost-Schur argument
\cite[Section~2.1]{DeLellisTopping} in the form needed here is as follows.
Solve $\Delta f=R-12$. The contracted Bianchi identity says
$\diver E=\tfrac14dR$, so integration by parts gives
\[
 \int_M(R-12)^2\dd=4\int_M\langle E,\nabla^2f\rangle\dd.
\]
Since $\Ric\ge0$, the integrated Bochner formula yields
$\norm{\nabla^2f}_2\le\norm{\Delta f}_2$. Cauchy--Schwarz proves the
second assertion in \eqref{eq:normalization}. Orthogonality of the trace
and trace-free parts gives
\begin{equation}\label{eq:ricci-error}
 \norm{\Ric-3g}_2^2
 =\norm E_2^2+\frac14\norm{R-12}_2^2\longrightarrow0.
\end{equation}
Moreover,
\[
 \int_MR^2\dd=144V+\norm{R-12}_2^2.
\]
Combining this identity with \eqref{eq:euler} and \eqref{eq:algebraic}
gives $V\ge v_0>0$ for all sufficiently large $j$.

It remains to control the diameter. Write $D=\diam(M,g)$ and choose
points realizing it. Let $A_1,A_2$ be the balls of radius $D/4$ about
these points. Bishop--Gromov comparison gives
$\Vol(A_i)\ge8^{-4}V$, by comparison with the ball of radius $2D$.
For any nonnegative continuous function $H$, the segment inequality is
\begin{equation}\label{eq:segment}
 \int_{A_1\times A_2}\int_{\gamma_{xy}}H\,ds\,d\mu(x)d\mu(y)
 \le CD\bigl(\Vol(A_1)+\Vol(A_2)\bigr)\int_MH\dd,
\end{equation}
where $\gamma_{xy}$ is the minimizing geodesic, unique for almost every
pair. This is the nonnegative-Ricci case of
\cite[Theorem~2.11]{CheegerColding}. To see the estimate directly, fix $x$
and use polar coordinates $y=\exp_x(r\theta)$ before cut time. On the
half-segment $r/2\le s\le r$, Jacobian comparison gives
\[
 J_x(r,\theta)\le(r/s)^3J_x(s,\theta)\le8J_x(s,\theta).
\]
Changing the order of integration bounds the integral over this half,
and over $y$, by $CD\int_MH\dd$. Integrating in $x\in A_1$ gives its
contribution to \eqref{eq:segment}. Reverse the geodesic and fix $y$
to treat the other half. The cut locus has measure zero.

Apply \eqref{eq:segment} to $H=|\Ric-3g|$. Some minimizing geodesic
between the two balls has length $\ell\ge D/2$ and satisfies
\begin{equation}\label{eq:good-segment}
 \int_0^\ell H(\gamma(s))\,ds\le\delta_jD,
 \qquad \delta_j=CV^{-1/2}\norm{\Ric-3g}_2\longrightarrow0.
\end{equation}
Use three parallel orthonormal normal fields along $\gamma$, each
multiplied by $\sin(\pi s/\ell)$, in the index form. Summing gives
\[
 \frac{3\pi^2}{2\ell}
 \ge\int_0^\ell\sin^2(\pi s/\ell)
             \Ric(\dot\gamma,\dot\gamma)\,ds
 \ge\frac{3\ell}{2}-\delta_jD.
\]
Because $D\le2\ell$, we obtain
$ (3/2-2\delta_j)\ell^2\le3\pi^2/2$.
Thus $D$ is uniformly bounded. Bishop comparison now gives the uniform
upper volume bound as well.
\end{proof}

\subsection{A support estimate and the Sobolev inequality}

We will need Sobolev inequalities both initially and along the flow.
The following elementary truncation lemma allows us to derive both
from estimates for functions with small support.

\begin{lemma}\label{lem:truncation}
Let $(M,g)$ be a closed four-manifold, $a>0$, and $P\ge0$ a smooth
function. Set
\[
 \E(f)=\int_M(a|\nabla f|^2+Pf^2)\dd.
\]
Suppose that for every Lipschitz function $f$, with
$m=\Vol\{f\ne0\}$, one has
\begin{equation}\label{eq:support}
 \norm f_2^2\le C m^{1/2}\E(f).
\end{equation}
Then $\norm f_4^2\le C'\E(f)$ for every such $f$, where $C'$ depends
only on $C$.
\end{lemma}

\begin{proof}
For $k\in\mathbb Z$ define
\[
 m_k=\Vol\{|f|>2^k\},\qquad
 f_k=\min\{(|f|-2^k)_+,2^k\}.
\]
The assumed estimate implies
\[
 2^{2k}m_{k+1}\le\norm{f_k}_2^2
 \le C m_k^{1/2}\E(f_k).
\]
The gradients of $f_k$ occupy disjoint layers up to null sets, and
$\sum_k f_k^2\le C_1f^2$ pointwise by a geometric series. Consequently
$\sum_k\E(f_k)\le C_2\E(f)$, with numerical constants. Put
$b_k=2^{4k}m_k$. The preceding estimate becomes
\[
 b_{k+1}\le16C b_k^{1/2}\E(f_k).
\]
The sum $B=\sum_kb_k$ is finite and comparable, with numerical constants,
to $\int_M|f|^4\dd$: integrate the geometric series
$\sum_{2^k<|f|}2^{4k}$ pointwise. Summing over $k$ and using $b_k^{1/2}\le B^{1/2}$
gives $B\le16CC_2B^{1/2}\E(f)$. Division by $B^{1/2}$ when $B>0$
proves the result; the case $B=0$ is immediate.
\end{proof}

\begin{proposition}\label{prop:initial-sobolev}
The normalized metrics of Proposition~\ref{prop:initial} satisfy
\begin{equation}\label{eq:initial-sobolev}
 \norm f_4^2\le S\int_M(|\nabla f|^2+f^2)\dd
\end{equation}
for all Lipschitz functions $f$, with a uniform constant $S$.
\end{proposition}

\begin{proof}
We establish \eqref{eq:support} with $a=P=1$.
Let $\bar f=V^{-1}\int_Mf\dd$. Polar integration of
$|f(x)-f(y)|$ along minimizing segments yields
\begin{equation}\label{eq:kernel}
 |f(x)-\bar f|
 \le\frac{CD^4}{V}\int_M\frac{|\nabla f(y)|}{d(x,y)^3}\,d\mu(y).
\end{equation}
Indeed, on a ray from $x$, the coefficient of $|\nabla f|$ at distance
$s$ is bounded using
\[
 \int_s^{\operatorname{cut}(\theta)}J_x(r,\theta)\,dr
 \le \frac{J_x(s,\theta)}{s^3}\int_s^D r^3\,dr
 \le\frac{D^4J_x(s,\theta)}{4s^3}.
\]
This proves \eqref{eq:kernel}; absolute continuity along almost every
ray justifies it for Lipschitz functions.

For any measurable set $A$ of volume $m>0$, comparison with Euclidean
polar measure gives
\[
 \sup_x\int_A d(x,y)^{-3}\,d\mu(y)\le C m^{1/4}.
\]
To verify this, split the integral at radius $\rho=m^{1/4}$. The inner
integral is at most $C\rho$ because $J_x(r,\theta)\le r^3$; the outer
integral is at most $m\rho^{-3}$. By symmetry the same bound holds
with $x$ and $y$ interchanged. The integral operator with this kernel,
restricted to $A$ in both variables, therefore has $L^2$ norm at most
$Cm^{1/4}$ by the Schur estimate.

Take $A=\{f\ne0\}$. The gradient vanishes almost everywhere off $A$,
and \eqref{eq:initial-geometry} bounds $D^4/V$. Taking $L^2(A)$ norms
in \eqref{eq:kernel} gives
\[
 \norm f_2\le Cm^{1/4}\norm{\nabla f}_2+\sqrt m\,|\bar f|
 \le Cm^{1/4}\norm{\nabla f}_2+\frac mV\norm f_2.
\]
If $m/V\le1/2$, absorb the last term and square. If $m/V>1/2$,
then $m\ge v_0/2$, and \eqref{eq:support} follows directly from the
$\int f^2$ term in the energy. Lemma~\ref{lem:truncation} completes
the proof.
\end{proof}

\section{A uniform smoothing interval}
\label{sec:smoothing}

Let $g_j(t)$ be the maximal smooth unnormalized Ricci flows with the
normalized initial metrics above:
\[
 \partial_tg_j=-2\Ric_{g_j},\qquad g_j(0)=g_j.
\]
Short-time existence and the curvature continuation criterion are
standard on closed manifolds \cite{HamiltonRF}. Initially each flow
has its own maximal time $\mathcal T_j>0$. Until a common lower bound
for these times is proved, every estimate refers only to existing
times. The scalar evolution and volume variation are
\begin{equation}\label{eq:scalar-volume}
 \partial_tR=\Delta R+2|\Ric|^2,\qquad
 \partial_td\mu=-R\,d\mu.
\end{equation}
Thus $R\ge0$ and $\Vol(M,g_j(t))\le V_0$ throughout smooth existence.
We first propagate a Sobolev inequality with a scalar-curvature weight,
then use smallness of the high-curvature part to remove that weight.

\subsection{Entropy supplies a weighted Sobolev inequality}

\begin{proposition}\label{prop:weighted}
There is a uniform $S_1$ such that, whenever
$0\le t<\min\{1,\mathcal T_j\}$,
\begin{equation}\label{eq:weighted}
 \norm f_4^2\le S_1\int_M\bigl(4|\nabla f|^2+(R+1)f^2\bigr)\dd
\end{equation}
for every Lipschitz function $f$ on $(M_j,g_j(t))$.
\end{proposition}

\begin{proof}
In dimension four, Perelman's entropy is
\begin{align*}
 \mu(g,\tau)=\inf_{\int f^2\dd=1}\bigg\{
 &\tau\int_M(4|\nabla f|^2+Rf^2)\dd
 -\int_Mf^2\log f^2\dd\\
 &-2\log(4\pi\tau)-4\bigg\},
\end{align*}
where the infimum is over smooth positive $f$. Its monotonicity gives
\begin{equation}\label{eq:perelman}
 \mu(g_j(t),\tau)\ge\mu(g_j(0),\tau+t),\qquad \tau>0
\end{equation}
\cite[Section~3]{Perelman}. Here $\tau+t$ is an entropy parameter, not
a later flow time, so \eqref{eq:perelman} imposes no additional lifespan
requirement.

At time zero, Jensen's inequality and \eqref{eq:initial-sobolev} give
\[
 \int_Mf^2\log f^2\dd\le\log\int_Mf^4\dd
 \le2\log\bigl(S(y+1)\bigr),\qquad
 y=\int_M|\nabla f|^2\dd,
\]
for $\int f^2\dd=1$. Since $R\ge0$,
\[
 \mu(g_j(0),\tau)\ge
 \inf_{y\ge0}\{4\tau y-2\log(\tau(y+1))\}-C\ge-C'
 \quad(0<\tau\le2).
\]
The last bound follows by writing $z=\tau(y+1)$: the expression in
braces is $4z-2\log z-4\tau$, uniformly bounded below in this range.
Consequently \eqref{eq:perelman} gives, at all times under consideration,
\begin{equation}\label{eq:log-sobolev}
 \int_Mf^2\log f^2\dd
 \le\tau\E_0(f)-2\log\tau+C,
 \qquad \E_0(f)=\int_M(4|\nabla f|^2+Rf^2)\dd,
\end{equation}
whenever $\norm f_2=1$ and $0<\tau\le1$.
Approximation of $|f|$ by positive smooth functions extends this to
Lipschitz $f$, with $0\log0=0$.

Let $m=\Vol\{f\ne0\}$. For a normalized $f$, Jensen's inequality on
its support gives $\int f^2\log f^2\dd\ge\log(1/m)$.
Choose $\tau=(1+\E_0(f))^{-1}$ in \eqref{eq:log-sobolev}. It follows
that
\[
 \log(1/m)\le2\log(1+\E_0(f))+C.
\]
Exponentiation, followed by homogeneity for arbitrary $f$, yields
\[
 \norm f_2^2\le Cm^{1/2}\left(\E_0(f)+\norm f_2^2\right).
\]
Apply Lemma~\ref{lem:truncation} with $a=4$ and $P=R+1$.
\end{proof}

\subsection{The high-curvature part stays small}

In orthonormal frames evolving by
$\partial_te_i=\Ric^\sharp e_i$, the curvature equation has the form
\begin{equation}\label{eq:curvature-flow}
 \partial_t\Rm=\Delta\Rm+\mathcal Q(\Rm),\qquad
 |\mathcal Q(\Rm)|\le C_q|\Rm|^2.
\end{equation}
Here time derivatives are derivatives of the moving-frame components;
the connection Laplacian has the sign fixed in Section~\ref{sec:intro}.
It follows that $u=|\Rm|$ satisfies
$(\partial_t-\Delta)u\le C_qu^2$ wherever $u>0$.

Choose the scalar function $K$ by
\begin{equation}\label{eq:threshold}
 K(0)=12C_0,\qquad K'=C_qK^2,
\end{equation}
where $C_0$ is from \eqref{eq:algebraic}. Fix $T\in(0,1/6)$ so small
that $K$ is bounded on $[0,T]$, and set
\[
 w_j(t)=\bigl(|\Rm_{g_j(t)}|-K(t)\bigr)_+.
\]
At time zero, \eqref{eq:algebraic} implies
\begin{equation}\label{eq:initial-tail}
 0\le w_j(0)\le C_0(R_{g_j}-12)_+,\qquad
 \norm{w_j(0)}_2\longrightarrow0.
\end{equation}

\begin{proposition}\label{prop:tail}
After discarding finitely many indices, there are uniform constants
$S_2,C$ such that, for $0\le t<\min\{T,\mathcal T_j\}$,
\begin{align}
 \norm f_4^2&\le S_2\int_M(|\nabla f|^2+f^2)\dd,
       \label{eq:flow-sobolev}\\
 \int_M w_j(t)^2\dd&\le e^{Ct}\int_M w_j(0)^2\,d\mu_{g_j},
       \label{eq:tail-propagation}\\
 \int_M|\Rm_{g_j(t)}|^2\dd&\le C.
       \label{eq:curvature-energy}
\end{align}
The Sobolev estimate holds for every Lipschitz function $f$.
In particular, the supremum of $\norm{w_j(t)}_2$ over these existing
times tends to zero.
\end{proposition}

\begin{proof}
We first work on any interval where $\norm w_2\le\eta$ for a small
fixed $\eta>0$. Since $R\le C|\Rm|\le C(K+w)$, the last term in
\eqref{eq:weighted} is bounded by a fixed multiple of $\int f^2$ plus
\[
 C\int_Mwf^2\dd\le C\eta\norm f_4^2.
\]
Choose $\eta$ so that this term is absorbed into the left side.
This gives \eqref{eq:flow-sobolev} with a fixed $S_2$; decreasing $\eta$
further does not change the chosen $S_2$.

On $\{u>K\}$, the equation for $K$ removes the bounded-background
reaction:
\[
 (\partial_t-\Delta)(u-K)\le C_q(u^2-K^2)
 =C_q(w^2+2Kw).
\]
Integration by parts and \eqref{eq:scalar-volume} therefore give
\begin{equation}\label{eq:tail-energy}
 \frac d{dt}\int_Mw^2\dd
 \le-2\int_M|\nabla w|^2\dd
       +2C_q\int_Mw^3\dd+4C_qK\int_Mw^2\dd.
\end{equation}
There is no adverse volume term, since $-\int Rw^2\dd\le0$.
The calculation is valid for the positive part: $K>0$, so $u$ is smooth
near $\{u\ge K\}$; the function $(u-K)_+^2$ is continuously
differentiable, and Lipschitz truncation justifies integration by parts.

H\"older's inequality and \eqref{eq:flow-sobolev} give
\[
 \int_Mw^3\dd\le\eta\norm w_4^2
 \le\eta S_2\int_M(|\nabla w|^2+w^2)\dd.
\]
Decrease $\eta$ so that $2C_q\eta S_2\le1$.
Then $Y(t)=\int w^2\dd$ satisfies $Y'\le CY$ as long as
$Y^{1/2}\le\eta$. By \eqref{eq:initial-tail}, for large $j$ we have
$Y(0)e^{CT}<\eta^2$. Gronwall's inequality and continuity exclude a
first time with $Y^{1/2}=\eta$ before $\min\{T,\mathcal T_j\}$.
This proves \eqref{eq:flow-sobolev} and \eqref{eq:tail-propagation}
on the entire existing interval. Finally,
$u\le K+w$ and $\Vol(M,g_j(t))\le V_0$ give
\eqref{eq:curvature-energy}.
\end{proof}

\subsection{Noncollapse and the common lifespan}

No common existence time has been used in this argument. We now use
its two outputs, uniform Sobolev control and a vanishing curvature tail,
to establish one.

\begin{lemma}\label{lem:noncollapse}
There is a uniform $v_*>0$ such that at every time covered by
Proposition~\ref{prop:tail},
\begin{equation}\label{eq:noncollapse}
 \Vol_{g_j(t)}B_{g_j(t)}(x,r)\ge v_*r^4
 \quad\text{for all }x\in M_j,\quad 0<r\le1.
\end{equation}
\end{lemma}

\begin{proof}
Write $V_x(r)=\Vol B(x,r)$ on one smooth slice. Insert into
\eqref{eq:flow-sobolev} the distance cutoff that is $1$ on $B(x,r/2)$,
vanishes outside $B(x,r)$, and has gradient at most $2/r$. Then
\[
 V_x(r/2)^{1/2}\le Cr^{-2}V_x(r).
\]
For $a(r)=V_x(r)/r^4$, this gives $a(r)\ge c\,a(r/2)^{1/2}$.
Iteration yields
\[
 a(r)\ge c^{2(1-2^{-k})}a(2^{-k}r)^{2^{-k}}.
\]
For this fixed smooth metric, $a(s)\to\omega_4$ as $s\downarrow0$,
where $\omega_4$ is the volume of the Euclidean unit ball in $\R^4$.
Letting $k\to\infty$ gives $a(r)\ge c^2$. No uniform scale for the
Euclidean asymptotic is required.
\end{proof}

\begin{proposition}\label{prop:smoothing}
For all sufficiently large $j$, the flow $g_j(t)$ exists on $[0,T]$.
There is a uniform constant $A$ such that
\begin{equation}\label{eq:curvature-bound}
 \max_{M_j}|\Rm_{g_j(t)}|\le\frac At,
 \qquad 0<t<T.
\end{equation}
\end{proposition}

\begin{proof}
Fix $A>2$ with $A/T\ge1$ and $A/T\ge4\max_{[0,T]}K$.
Suppose that $t\max_Mu(t)$ first reaches $A$ at an existing time $t<T$.
Set $Q=\max_Mu(t)=A/t$. On $[t-Q^{-1},t]$ the curvature is at most
$2Q$, because $t-Q^{-1}\ge t/2$ and this is a first crossing.
Parabolic rescaling gives a smooth closed flow on $[0,1]$ whose
curvature norm is at most $2$. Shi's first derivative estimate
\cite{Shi} gives
\[
 |\nabla\Rm|(t)\le CQ^{3/2}.
\]
This estimate requires no initial derivative or injectivity-radius
bound; one may use the local formulation in
\cite[Theorem~3.27]{MorganTian} with the entire closed manifold as its
domain.

Let $x$ realize the curvature maximum. Choose a fixed $c\in(0,1]$
small enough that on $B_{g_j(t)}(x,cQ^{-1/2})$ the derivative estimate
implies $u\ge3Q/4$. Since $Q\ge4K(t)$, we have $w\ge Q/2$ there.
Also $Q\ge1$, so this radius is at most $1$. Lemma~\ref{lem:noncollapse}
now gives
\[
 \int_Mw(t)^2\dd\ge\frac{Q^2}{4}\,v_*c^4Q^{-2}
 =\frac{v_*c^4}{4}.
\]
This contradicts the uniform convergence of the tail norm to zero in
Proposition~\ref{prop:tail}. Thus, for large $j$, no crossing occurs,
and \eqref{eq:curvature-bound} holds throughout existing times below $T$.

If $\mathcal T_j\le T$, the same bound controls curvature on
$[\mathcal T_j/2,\mathcal T_j)$. Curvature is already bounded on the
earlier compact interval by smoothness of the individual flow. The
curvature continuation criterion contradicts maximality. Hence
$\mathcal T_j>T$.
\end{proof}

\section{The Einstein and sectional-curvature defects}
\label{sec:defects}

The common interval now permits compactness at positive times. To
identify the eventual limit, we must retain information from time zero:
the Einstein error tends to zero, and the integral negative sectional
curvature remains of order $t$.

\begin{proposition}\label{prop:einstein-defect}
Set $\lambda(t)=3/(1-6t)$. Along the flows above,
\begin{equation}\label{eq:einstein-defect}
 \sup_{0\le t<T}
 \norm{\Ric_{g_j(t)}-\lambda(t)g_j(t)}_{L^2(g_j(t))}
 \longrightarrow0.
\end{equation}
\end{proposition}

\begin{proof}
Write $L=\Ric-\lambda g$. In moving orthonormal frames, define the
curvature action on a symmetric tensor $B$ by
\[
 (\Rm(B))_{ij}=\sum_{k,l}R_{ikjl}B_{kl};
 \qquad \Rm(g)=\Ric.
\]
The Ricci evolution in these frames is
\[
 \partial_t\Ric=\Delta\Ric+2\Rm(\Ric).
\]
The frame motion cancels the $-2\Ric^2$ term in the fixed-frame formula;
the components of $g$ in the moving frames are constant. Since
$\lambda'=2\lambda^2$, subtraction yields
\begin{equation}\label{eq:L-flow}
 \partial_tL=\Delta L+2\Rm(L)+2\lambda L.
\end{equation}
Using $R\ge0$, $|\Rm|\le K+w$, and boundedness of $K$ and $\lambda$
on $[0,T]$, we obtain
\[
 \frac d{dt}\int_M|L|^2\dd
 \le-2\int_M|\nabla L|^2\dd
       +C\int_M(1+w)|L|^2\dd.
\]
By H\"older, \eqref{eq:flow-sobolev}, and Kato's inequality,
\[
 \int_Mw|L|^2\dd
 \le\norm w_2\norm{|L|}_4^2
 \le S_2\norm w_2\int_M(|\nabla L|^2+|L|^2)\dd.
\]
The coefficient tends to zero uniformly in time, by
\eqref{eq:tail-propagation}. Absorb the gradient term for large $j$
and apply Gronwall's inequality. The initial norm tends to zero by
\eqref{eq:ricci-error}, proving \eqref{eq:einstein-defect}.
\end{proof}

We next measure failure of nonnegative sectional curvature by a scalar
function of the curvature tensor. On the Euclidean space of algebraic
curvature tensors on $\R^4$, define
\begin{equation}\label{eq:F-definition}
 F(B)=\max\left\{0,
 \sup_{\substack{|X|=|Y|=1\\\langle X,Y\rangle=0}}
              -B(X,Y,X,Y)\right\}.
\end{equation}
Thus $F(\Rm)$ is the magnitude of the most negative sectional curvature,
or zero if all sectional curvatures are nonnegative.

\begin{proposition}\label{prop:sectional-defect}
There is a uniform constant $C$ such that
\begin{equation}\label{eq:sectional-defect}
 \int_{M_j}F(\Rm_{g_j(t)})\,d\mu_{g_j(t)}\le Ct,
 \qquad 0\le t<T.
\end{equation}
\end{proposition}

\begin{proof}
The function $F$ is nonnegative, convex, orthogonally invariant, and
Lipschitz: it is a supremum of a bounded family of linear functionals
and the zero function. Let $F_\delta$ be its convolution with a
nonnegative smooth radial mollifier supported in a ball of radius
$\delta$ in the algebraic-curvature tensor space. Then $F_\delta$ is
smooth, nonnegative, convex, and orthogonally invariant, with
\[
 |DF_\delta|\le C,\qquad |F_\delta-F|\le C\delta.
\]
Orthogonal invariance makes this construction independent of the
orthonormal frame and compatible with the metric connection.

Apply the chain rule to \eqref{eq:curvature-flow} in moving
orthonormal frames. Convexity gives a favorable Hessian term:
\begin{align*}
 (\partial_t-\Delta)F_\delta(\Rm)
 &=DF_\delta(\Rm)[\mathcal Q(\Rm)]
   -\sum_aD^2F_\delta(\Rm)
           [\nabla_a\Rm,\nabla_a\Rm]\\
 &\le C|\Rm|^2.
\end{align*}
After integration, the volume variation contributes
$-\int R F_\delta(\Rm)\dd\le0$. Closedness and
\eqref{eq:curvature-energy} therefore imply
\[
 \int_MF_\delta(\Rm(t))\dd
 \le\int_MF_\delta(\Rm(0))\,d\mu_{g_j}+Ct
 \le C\delta V_0+Ct.
\]
The last step uses $F(\Rm(0))=0$. Letting $\delta\downarrow0$ proves
\eqref{eq:sectional-defect}.
\end{proof}

\section{The compact Einstein limit and the gap}
\label{sec:limit}

We first extract the geometric conclusion that does not use
Classification Assumption~\ref{ass:classification}.

\begin{proposition}\label{prop:limit}
For any sequence \eqref{eq:sequence}, normalize the metrics as in
Proposition~\ref{prop:initial} and discard finitely many terms.
There are a subsequence, a connected simply connected closed smooth
four-manifold $N$, a metric $h$ on $N$, and diffeomorphisms
$\Phi_j:N\to M_j$ such that
\[
 \Ric_h=3h,\qquad \sec_h\ge0,
\]
and
\begin{equation}\label{eq:flow-limit}
 \Phi_j^*g_j(t)\longrightarrow (1-6t)h
\end{equation}
smoothly on $N\times I$ for every compact interval $I\subset(0,T)$.
\end{proposition}

\begin{proof}
Fix $t_*=T/2$. By \eqref{eq:curvature-bound}, the flows have uniform
global curvature bounds on every compact time interval in $(0,T)$.
At $t_*$, the local volume lower bound \eqref{eq:noncollapse} and the
total volume upper bound $V_0$ also give a diameter bound. Indeed, points
spaced a fixed distance apart along a minimizing geodesic have pairwise
disjoint balls of a fixed smaller radius; each such ball has uniformly
positive volume. Their number is bounded by $V_0$. Write the resulting
bound as $\diam(M_j,g_j(t_*))\le D_*$.

Absolute sectional curvature bounds and the ball-volume lower bound
give a uniform positive injectivity radius at $t_*$, by the injectivity
radius estimate of Cheeger--Gromov--Taylor
\cite[Theorem~4.7]{CheegerGromovTaylor}. Hamilton's compactness theorem
\cite{HamiltonCompactness} now gives, after passing to a subsequence,
smooth pointed convergence to a connected Ricci flow $(N,h(t))$ on
$(0,T)$, with $h(t_*)$ complete. One can apply the theorem on nested
compact time intervals containing $t_*$ and take a diagonal subsequence.
The curvature bounds used here are global in space; hence the complete
limit version of compactness applies, without the completeness issue
that can arise from purely local bounds
\cite[Theorem~1.1 and subsequent discussion]{ToppingCompactness}.
The convergence maps are time-independent embeddings defined on an
exhaustion of $N$.

We explain why the limit is compact and the embeddings are eventually
diffeomorphisms. If the complete connected slice $(N,h(t_*))$ were
noncompact, it would contain a compact closed ball $\overline B(p,L)$
with a point $q$ at distance $L$, for any $L>0$. Choose $L>2D_*$.
For large $j$, a convergence embedding is defined on a neighborhood of
this ball, and its pulled-back metric bounds lengths on the ball from
below by half their $h(t_*)$ values. Any curve in $M_j$ joining the
images of $p$ and $q$ must have length at least $L/2$: if it leaves the
image of the ball, its initial portion must first reach that image's
boundary and already has this length. If it stays inside, the same
bound follows by pullback. This contradicts the diameter bound.
Thus $N$ is compact. Eventually the convergence embeddings are defined
on all of $N$; their images are open, by the inverse function theorem,
and closed, by compactness. Since $M_j$ is connected, each is a
diffeomorphism. It follows also that $N$ is simply connected.

Smooth convergence on this compact manifold and
\eqref{eq:einstein-defect} give
\[
 \Ric_{h(t)}=\lambda(t)h(t),\qquad 0<t<T.
\]
The Ricci-flow equation therefore integrates to
\[
 h(t)=(1-6t)h,\qquad
 h=\frac{1}{1-6t_*}h(t_*),\qquad \Ric_h=3h.
\]
For each fixed $t>0$, \eqref{eq:sectional-defect} also passes to the
limit. The function $F$ is positively homogeneous of degree one.
Under metric scaling by a positive constant $c$, its value on curvature
scales by $c^{-1}$ and the four-dimensional volume form scales by $c^2$.
Consequently
\begin{equation}\label{eq:final-scaling}
 (1-6t)\int_NF(\Rm_h)\,d\mu_h\le Ct,
 \qquad 0<t<T.
\end{equation}
Let $t\downarrow0$. The nonnegative continuous integrand must vanish,
so $\sec_h\ge0$. This proves the proposition.
\end{proof}

\begin{proof}[Proof of Theorem~\ref{thm:gap}]
If no universal positive gap existed, then for every positive integer
$j$ there would be a connected simply connected closed four-manifold
$(M_j,g_j)$ with $\sec_{g_j}\ge0$ and $\mathcal D(g_j)<1/j$, whose
smooth manifold is outside the three stated diffeomorphism types.
Proposition~\ref{prop:limit} gives a subsequence diffeomorphic to a
closed simply connected $(N,h)$ with $\Ric_h=3h$ and $\sec_h\ge0$.
Classification Assumption~\ref{ass:classification} applies exactly
to this metric. Since $N$ is already simply connected, it is one of
the three model manifolds, a contradiction.
\end{proof}

\begin{remark}
Proposition~\ref{prop:limit} is independent of the classification
input. It asserts smooth convergence only after positive-time
regularization. In particular, the recovery of $\sec_h\ge0$ in
\eqref{eq:final-scaling} requires neither convergence of the original
metrics at time zero nor preservation of nonnegative sectional
curvature by Ricci flow.
\end{remark}

\bibliographystyle{abbrv}
\bibliography{references}
\end{document}
